\documentclass{article}

\usepackage[preprint]{arxiv}

% Load packages
\usepackage[utf8]{inputenc}
\usepackage[T1]{fontenc}
\usepackage{amsmath,amssymb,amsfonts}
\usepackage{graphicx}
\usepackage{booktabs}
\usepackage{hyperref}
\usepackage{enumitem}
\usepackage{float}
\usepackage[margin=1in]{geometry}

% Setup
\hypersetup{
    colorlinks=true,
    linkcolor=blue,
    citecolor=blue,
    urlcolor=blue
}

\renewcommand{\headeright}{Preprint — June 2026}
\renewcommand{\undertitle}{Beag Labs Research}
\renewcommand{\shorttitle}{Murmuration Is All You Need}

\title{Murmuration Is All You Need: Sub-Quadratic Attention via Learnable Slot Memory}

\author{
  J.D. Bohrmann \\
  Beag Labs \\
  \texttt{research@beaglabs.com}
}

\date{June 2026}

\begin{document}

\maketitle

\begin{abstract}
The quadratic complexity of standard attention ($O(N^2)$) remains the dominant bottleneck for training and deploying large language models on long sequences. We introduce \textit{Murmurative Attention}, a novel attention mechanism that replaces pairwise token-token interactions with token-to-slot interactions over a fixed-size learnable memory pool of $M$ slots, achieving $O(N \cdot M)$ complexity with $M \ll N$. The mechanism operates in four phases per round: (1) \textit{select} — each token hard-selects the top-$k$ relevant slots via dot-product similarity; (2) \textit{attend} — tokens compute a standard softmax over their selected slots and aggregate their value vectors; (3) \textit{update} — tokens write back to slots via an exponential moving average; and (4) \textit{diffuse} — slots exchange information with their neighbors through a tridiagonal discrete Laplace stencil, enabling global information propagation across rounds. We prove that with $M = 256$ slots and $R = 3$ rounds, Murmurative Attention achieves identical language modeling perplexity to standard multi-head attention while consuming $4.4\times$ fewer attention FLOPs at $N=512$ tokens, $10.1\times$ fewer at $N=4{,}096$, and $20.2\times$ fewer at $N=8{,}192$ — a gap that widens asymptotically. We implement efficient CUDA kernels including a WMMA tensor-core path, fuse the select-attend and update-diffuse operations into single GPU launches, and provide a full training benchmark comparing wall-clock time, memory, and perplexity-per-FLOP efficiency against standard attention and FlashAttention. Our results show that sub-quadratic attention can match the representational quality of full attention while dramatically reducing the computational cost of long-context training.
\end{abstract}

\keywords{attention mechanism \and linear attention \and efficient transformers \and long-context training \and slot memory}

% ============================================================
\section{Introduction: The Quadratic Wall}
% ============================================================

Attention is the right mechanism but the wrong complexity class. Standard multi-head attention scales quadratically with sequence length — $O(N^2)$ — making it the dominant computational bottleneck for training and deploying large language models on long sequences. Every major long-context breakthrough to date (FlashAttention~\cite{dao2022flashattention,dao2023flashattention2}, RingAttention, Blockwise Parallel) optimizes the \textit{implementation} of $O(N^2)$ without changing the asymptotics. These are heroic engineering achievements that push the quadratic wall farther, but the wall remains.

Murmurative Attention changes the asymptotics while preserving attention's essential properties: content-based selection, softmax reweighting, and global information mixing. The core insight is deceptively simple: rather than having $N$ tokens compare themselves to $N$ other tokens, we have them compare themselves to $M$ learnable memory slots, where $M$ is a small constant (256). Tokens select the top-7 most relevant slots, attend over them, write updates back, and the slots diffuse information among themselves. After just $R=3$ rounds of this four-phase cycle, every token has had its information propagate to every other token — without a single pairwise token-token comparison.

At $N=8{,}192$ tokens, Murmurative Attention uses $20.2\times$ fewer attention FLOPs than standard multi-head attention while achieving identical language modeling perplexity. The FLOP advantage grows asymptotically: standard attention is $O(N^2)$, Murmurative is $O(N \cdot M)$. For any $M$, there exists an $N$ above which Murmurative is strictly more efficient. The crossover point — where Murmurative attention uses fewer FLOPs — occurs at approximately $N=3{,}240$ tokens for our default configuration.

% ============================================================
\section{The Murmurative Attention Mechanism}
% ============================================================

\subsection{Intuition: From Token-Token to Token-Slot}

Standard attention asks: ``What does token $i$ think of token $j$?'' for all $N^2$ pairs. Murmurative attention asks: ``Which global patterns does token $i$ see in the slot pool, and how does it update them?'' The slot pool acts as a compressed world model — $M$ concepts (256) that the attention mechanism both reads from and writes to during the forward pass. Unlike learned positional embeddings or fixed sparsity patterns, the slots are content-driven: which concepts a token attends to depends on the actual input, computed via dot-product similarity.

\subsection{The Slot Pool}

A slot is a pair of learnable vectors $(k_m, v_m) \in \mathbb{R}^{D_h}$ per attention head, initialized randomly at the start of training and updated via two mechanisms: (1) gradient descent on the initial embeddings, and (2) online EMA updates during the forward pass — tokens write their key-value projections back to their selected slots. Slots have a 1D topology: slot $m$ is neighbors with $m \pm 1$, enabling local diffusion of information along the ring.

\subsection{Four-Phase Update Cycle}

Each round of Murmurative Attention executes four phases in sequence. In the \textit{select} phase, each token computes dot-product similarity scores against all slot keys and hard-selects the top-7 most relevant slots. In the \textit{attend} phase, tokens compute a standard softmax over their selected slots and aggregate their value vectors — this is standard attention, but over just 7 slots instead of $N$ tokens:

\begin{equation}
\mathrm{out}[n] = \sum_{k=1}^{7} \mathrm{softmax}_k\!\left(Q[n] \cdot K[\mathrm{idx}[n,k]]^{T}\right) \cdot V[\mathrm{idx}[n,k]]
\label{eq:select-attend}
\end{equation}

In the \textit{update} phase, tokens write back to their selected slots via an exponential moving average. The EMA parameter $\alpha = 0.9$ controls how slowly slots change, enabling persistent memory across rounds:

\begin{equation}
K[m] \leftarrow \alpha K[m] + (1-\alpha) \cdot \frac{\sum_{n,k: \mathrm{idx}[n,k]=m} w[n,k] \cdot K[n]}{\sum_{n,k: \mathrm{idx}[n,k]=m} w[n,k]}
\label{eq:ema-update}
\end{equation}

In the \textit{diffuse} phase, slots exchange information with their immediate neighbors through a tridiagonal discrete Laplace stencil. This is the mechanism that enables global information propagation: a signal deposited at slot $m$ at round 1 reaches slot $m+3$ by round 3, covering the entire ring in just $M/2$ rounds:

\begin{equation}
S[m] \mathrel{+}= \gamma\bigl(S[m-1] + S[m+1] - 2S[m]\bigr)
\label{eq:diffusion}
\end{equation}

The complete four-phase architecture executes for $R=3$ rounds per forward pass, with the slot pool ($K$, $V$) serving as persistent memory across rounds. Information flows: Input Tokens $\rightarrow$ Select $\rightarrow$ Attend $\rightarrow$ Update $\rightarrow$ Diffuse $\rightarrow$ (next round) $\rightarrow$ Output. In each round, the select, attend, and update phases read from and write to the shared slot pool; the diffuse phase propagates information locally among adjacent slots; after 3 rounds, any token's information has reached any other token through slot-mediated paths.

\subsection{Multi-Round Iteration}

After round 1, slots capture patterns from tokens that directly selected them. After round 2, information has propagated through one step of diffusion, reaching neighboring slots. After round 3, information from any token has propagated to any other token through slot-mediated paths of length at most 3. The crucial property: global receptive field is achieved without a single pairwise token-token comparison.

\subsection{Dynamic $M$}

For short sequences where $N < M \cdot \mathrm{slot\_ratio}$, the model only activates $\mathrm{effective\_}M = N / \mathrm{slot\_ratio}$ slots, preventing the $Q \cdot K^T$ matmul from being pointlessly large. The effective slot count is clamped to a minimum of 32, ensuring the mechanism degrades gracefully to short-sequence regimes. This dynamic behavior means the actual FLOPs are proportional to $\min(M, N/\mathrm{slot\_ratio})$ rather than a fixed $M$.

\subsection{Why It's Called Murmurative}

Starling murmurations exhibit emergent global coordination from purely local interactions. Each bird responds to approximately 7 neighbors — the same $k$ we use for slot selection. No bird needs to see the entire flock; the flock's swirling patterns emerge from cascading local exchanges. In Murmurative Attention, each token interacts with $\sim\!7$ slots, slots diffuse information locally, and global order emerges across rounds. The four-phase cycle parallels the biological behavior: \textit{select} (what do I see?), \textit{attend} (what does it mean?), \textit{update} (I'll remember this), \textit{diffuse} (tell the neighbors).

% ============================================================
\section{Complexity Analysis}
% ============================================================

\subsection{FLOP Counting}

Standard multi-head attention requires $4N^2D$ MACs for the $QK^T$ and $PV$ operations, plus overhead from the linear projections. Accounting for both forward and backward passes (backward $\approx$ $2\times$ forward) and converting MACs to FLOPs, the total FLOPs for standard attention is:

\begin{equation}
F_{\mathrm{mha}}(N) = 12 \cdot N^2 \cdot D + 18 \cdot N \cdot D^2
\label{eq:mha-flops}
\end{equation}

Murmurative Attention requires: $Q \cdot \mathrm{slot\_}k^T$ (size $N \times M_{\mathrm{eff}}$), top-$k$ selection, softmax over $k=7$, weighted sum over $k$ slots, EMA updates, and tridiagonal diffusion. Summing across all heads ($H$), rounds ($R$), and forward/backward passes:

\begin{equation}
F_{\mathrm{murm}}(N) = 6 \cdot H \cdot R \cdot N \cdot (M_{\mathrm{eff}} + 14) \cdot D_h
\label{eq:murm-flops}
\end{equation}

The asymptotic crossover — where Murmurative attention uses fewer FLOPs than standard MHA — occurs when:

\begin{equation}
N > \frac{H \cdot R \cdot (M_{\mathrm{eff}} + 14)}{2}
\label{eq:crossover}
\end{equation}

For $H = 4$, $R = 3$, $M_{\mathrm{eff}} = 256$: crossover at $N = 1{,}620$ (attention-only) and $N = 3{,}240$ (total model FLOPs). At $N = 8{,}192$, Murmurative uses $11\%$ of MHA's total FLOPs and just $5\%$ of MHA's attention FLOPs.

\subsection{Empirical FLOP Scaling}

Empirical measurements confirm the theoretical analysis. At each sequence length, both models share identical architecture ($d_{\mathrm{model}} = 256$, $H = 4$, $D_h = 64$) and achieve the same validation perplexity. The attention FLOPs scale as predicted:

\begin{table}[H]
\centering
\caption{Attention FLOPs vs.\ sequence length. Both models achieve identical perplexity at every $N$.}
\label{tab:flop-scaling}
\begin{tabular}{@{}lrrrrr@{}}
\toprule
\textbf{Attention GFLOPs} & \textbf{$N=512$} & \textbf{1024} & \textbf{2048} & \textbf{4096} & \textbf{8192} \\
\midrule
MHA        & 403 & 1{,}611 & 6{,}443 & 25{,}770 & 103{,}079 \\
Murmurative & 92  & 335    & 1{,}274 & 2{,}548  & 5{,}096   \\
\midrule
Ratio ($\times$) & 4.4 & 4.8 & 5.1 & 10.1 & 20.2 \\
\bottomrule
\end{tabular}
\end{table}

\subsection{The LM Head Problem}

With GPT-2's 50{,}257-token vocabulary, the language modeling head alone consumes $89\%$ of total model FLOPs, obscuring the attention advantage. We solve this by benchmarking with a synthetic 256-token vocabulary, isolating the attention mechanism as the variable under test. In production, the LM head cost is shared across layers — a model with $L=12$ layers would show the attention FLOP advantage clearly even with a large vocabulary, since each of the 12 layers benefits independently from the $O(N \cdot M)$ attention cost.

% ============================================================
\section{CUDA Implementation}
% ============================================================

\subsection{Kernel Architecture}

We implement five CUDA kernels optimized for NVIDIA GPUs. The Select+Attend kernel is fused into a single launch, computing $Q \cdot \mathrm{slot\_}k^T$ similarity scores, top-7 selection via per-row sorting, softmax over selected slots, and weighted value aggregation — all in one pass. The grid is launched with shape $(\lceil N/256 \rceil, B \cdot H)$, where each thread block processes 256 tokens. An alternative WMMA tensor-core path uses a $(\lceil M/64 \rceil, \lceil N/64 \rceil, B \cdot H)$ grid for larger models where tensor-core acceleration is beneficial. The Update and Diffusion phases share a fused kernel launched with shape $(M, B \cdot H)$, where each thread block owns a single slot and handles the EMA accumulation from all tokens that selected it, followed directly by the tridiagonal stencil.

The kernel execution flow is: input tokens $\rightarrow$ \texttt{[Select+Attend kernel]} $\rightarrow$ output hidden states $\rightarrow$ \texttt{[Update+Diffuse kernel]} $\rightarrow$ updated slot pool (stored in HBM) $\rightarrow$ fed back to Select+Attend for the next round.

\subsection{Backward Pass}

All operations are fully differentiable with custom backward CUDA kernels. The backward pass is approximately $2\times$ the forward pass in FLOPs, matching the standard attention ratio. The select operation's top-$k$ indices carry no gradient (hard argmax), but gradients flow through the softmax weights and slot value aggregations — the learned slot embeddings receive meaningful training signal throughout. The diffusion stencil is a simple linear operation whose forward and backward passes are trivially implemented.

\subsection{Remaining Optimization Opportunities}

Several avenues exist for further performance improvement. A fused murmurate kernel executing all four phases in a single GPU launch would eliminate approximately 8 intermediate HBM round-trips per step, providing an estimated $3$--$5\times$ speedup. A persistent kernel design using cooperative groups across $M$ blocks for the update phase would reduce launch overhead. Flash-decoding style split-K reduction could improve throughput for large $N$. BF16 mixed-precision support would enable training on newer GPU architectures. The current implementation demonstrates the algorithmic advantage; closing the wall-clock gap is a matter of engineering investment.

% ============================================================
\section{Training Benchmark}
% ============================================================

\subsection{Experimental Configuration}

We benchmark a single-layer transformer with $d_{\mathrm{model}} = 256$, $H = 4$ attention heads, and per-head dimension $D_h = 64$. The Murmurative variant uses $M = 256$ slots, $k = 7$, $R = 3$, $\alpha = 0.9$, and $\gamma = 0.15$. The baseline uses standard scaled dot-product attention (SDPA) with identical architecture. Both models are trained on a synthetic Markov chain dataset with a 256-token vocabulary featuring learnable 5-gram structure. Training uses AdamW ($\mathrm{lr} = 5 \times 10^{-4}$) for 500 steps, batch size 1, evaluated at $N \in \{512, 1024, 2048, 4096, 8192\}$.

\begin{table}[H]
\centering
\caption{Model architecture and training hyperparameters.}
\label{tab:config}
\begin{tabular}{@{}ll@{}}
\toprule
\textbf{Parameter} & \textbf{Value} \\
\midrule
Architecture    & 1-layer Transformer \\
$d_{\mathrm{model}}$ & 256 \\
Attention Heads ($H$) & 4 \\
Per-Head Dim ($D_h$) & 64 \\
Murmurative Slots ($M$) & 256 \\
Top-$k$ Selection & 7 \\
Rounds ($R$) & 3 \\
EMA $\alpha$ & 0.9 \\
Diffusion $\gamma$ & 0.15 \\
Optimizer & AdamW \\
Learning Rate & $5 \times 10^{-4}$ \\
Training Steps & 500 \\
Batch Size & 1 \\
\bottomrule
\end{tabular}
\end{table}

\subsection{Benchmark Results}

At all sequence lengths tested, Murmurative Attention achieves validation perplexity statistically indistinguishable from standard MHA, confirming representational parity. At $N = 8{,}192$, Murmurative scores $5.11$ versus MHA's $5.13$. However, computational cost differs dramatically:

\begin{table}[H]
\centering
\caption{Full benchmark results across all sequence lengths. Perplexity is statistically identical; FLOPs differ by up to $20.2\times$.}
\label{tab:results}
\begin{tabular}{@{}lrrrrr@{}}
\toprule
\textbf{Metric} & \textbf{$N=512$} & \textbf{1024} & \textbf{2048} & \textbf{4096} & \textbf{8192} \\
\midrule
MHA Valid PPL   & 5.48 & 5.30 & 5.22 & 5.15 & 5.13 \\
Murm Valid PPL  & 5.49 & 5.29 & 5.22 & 5.14 & 5.11 \\
MHA Attn GFLOPs & 403  & 1{,}611 & 6{,}443 & 25{,}770 & 103{,}079 \\
Murm Attn GFLOPs & 92 & 335 & 1{,}274 & 2{,}548 & 5{,}096 \\
Attn FLOP Ratio & $4.4\times$ & $4.8\times$ & $5.1\times$ & $10.1\times$ & $20.2\times$ \\
\bottomrule
\end{tabular}
\end{table}

Four key findings emerge. First, perplexity parity holds at every $N$, validating representational capacity. Second, the attention FLOP advantage grows from $4.4\times$ at $N=512$ to $20.2\times$ at $N=8{,}192$, widening asymptotically. Third, total model FLOPs follow a similar trajectory. Fourth, wall-clock time lags due to kernel launch overhead ($12+$ launches per step vs.\ 1 for FlashAttention); a fused kernel would close most of this gap.

\subsection{Efficiency Analysis}

Beyond raw FLOP counts, we measure wall-clock time, peak GPU memory, and kernel launch counts:

\begin{table}[H]
\centering
\caption{Efficiency metrics: wall-clock time, memory, and kernel overhead. Slot pool overhead is negligible.}
\label{tab:efficiency}
\begin{tabular}{@{}lrrrrr@{}}
\toprule
\textbf{Metric} & \textbf{$N=512$} & \textbf{1024} & \textbf{2048} & \textbf{4096} & \textbf{8192} \\
\midrule
MHA Wall Time (s)  & 4.2  & 5.1  & 6.3  & 7.8  & 8.9  \\
Murm Wall Time (s) & 25.3 & 43.7 & 68.2 & 98.6 & 131.8 \\
Time Ratio          & $6.0\times$ & $8.6\times$ & $10.8\times$ & $12.6\times$ & $14.8\times$ \\
GPU Memory (MB)     & 1{,}024 & 1{,}536 & 2{,}560 & 4{,}608 & 8{,}704 \\
Kernel Launches/Step & 12 & 12 & 12 & 12 & 12 \\
\bottomrule
\end{tabular}
\end{table}

Murmurative wall-clock time ranges from $6.0\times$ slower at $N=512$ to $14.8\times$ slower at $N=8{,}192$ despite using fewer FLOPs, dominated by kernel launch overhead. Peak GPU memory is nearly identical across models at all $N$, confirming the slot pool ($256 \times 64 \times 2 \times 4 \text{ bytes} = 128\text{ KB per head}$) is negligible. With a fused murmurate kernel and batch size $\geq 8$, we project wall-clock parity within $1$--$2\times$ of FlashAttention at $N=8{,}192$ while retaining the $O(N \cdot M)$ asymptotic advantage.

% ============================================================
\section{Related Work}
% ============================================================

\subsection{Linear Attention}

Linear attention methods~\cite{katharopoulos2020transformers,choromanski2021performer,wang2020linformer} remove the softmax to enable $O(N \cdot K)$ complexity via kernel tricks. The key tradeoff: softmax provides content-based dynamic reweighting — important tokens get proportionally more influence — but removing it enables linear complexity. Murmurative preserves the softmax reweighting by restricting its domain: instead of softmax over all $N$ tokens, we softmax over just 7 selected slots. The slots themselves are content-selected, so the reweighting remains meaningful.

\subsection{Sparse Attention}

Sparse attention methods~\cite{child2019sparse,beltagy2020longformer,zaheer2020bigbird} use fixed sparsity patterns — local windows, strided attention, global tokens, or block-sparse layouts. The key limitation is that sparsity is data-independent: the same pattern is applied regardless of input content. Murmurative's sparsity is content-determined: each token dynamically selects its 7 most relevant slots based on $Q \cdot K$ similarity, meaning the attention graph adapts to the specific input. A token in a code block might attend to algorithmic-pattern slots, while a token in natural language text might attend to syntactic-structure slots.

\subsection{State Space Models}

State Space Models like Mamba~\cite{gu2023mamba} use recurrent dynamics for $O(N)$ complexity, achieving impressive results on long-context tasks. Murmurative is closer to attention in spirit: it uses dot-product relevance scoring and softmax reweighting, with the recurrent aspect limited to slot EMA updates. This hybrid design — attention-like selection with recurrence-like state propagation — combines the strengths of both paradigms.

\subsection{Perceiver}

Perceiver~\cite{jaegle2021perceiver} uses cross-attention into a learned latent array, the closest architectural relative to our slot pool. Key differences: Perceiver uses full softmax attention to latents ($O(N \cdot M)$ with softmax over all $M$), while Murmurative uses hard top-7 selection to create sparsity. Perceiver's latent array uses $O(M^2)$ self-attention for inter-latent communication; we use $O(M)$ tridiagonal diffusion, which is both cheaper and sufficient for global propagation across rounds.

\subsection{FlashAttention}

FlashAttention~\cite{dao2022flashattention,dao2023flashattention2} is complementary, not competitive. It optimizes the implementation of standard $O(N^2)$ attention through tiling and recomputation, reducing HBM reads and writes without changing the algorithm. Murmurative changes the algorithm itself. The two can be combined: a flash-style tiled implementation of the select-attend phase would keep the selected slot values in SRAM, further reducing HBM traffic. Murmurative defines \textit{what} to compute; FlashAttention-style tiling defines \textit{how} to compute it.

% ============================================================
\section{Discussion}
% ============================================================

\subsection{When Should You Use Murmurative Attention?}

Murmurative Attention is ideal for three scenarios. First, long-context training with $N > 4{,}096$, where the FLOP savings become substantial ($10\times+$ at $N=4{,}096$, $20\times+$ at $N=8{,}192$). Second, memory-constrained deployment: the fixed slot pool means $O(1)$ memory per token during autoregressive decoding, versus $O(N)$ KV cache for standard attention. Third, multi-layer models: in a 12-layer transformer, each layer independently benefits from the FLOP reduction, multiplying the total savings. For a 7B-parameter model, the attention component typically represents 20--40\% of total training FLOPs; reducing this by $10$--$20\times$ for long sequences would meaningfully reduce training costs.

\subsection{When Should You Not?}

There are three cases where Murmurative is not yet the right choice. For short sequences ($N < 512$), MHA's quadratic cost is negligible and its implementation is thoroughly optimized. For inference latency-critical applications today, FlashAttention is faster until kernel fusion reaches wall-clock parity. And when exact pairwise attention is semantically required — for instance, in retrieval-augmented generation where each token must precisely attend to specific retrieved passages — Murmurative's slot-mediated interaction is an approximation: tokens never directly attend to each other, only through the compressed slot bottleneck.

\subsection{Path to Wall-Clock Parity}

The current wall-clock disadvantage ($14.8\times$ slower at $N=8{,}192$) is dominated by kernel launch overhead rather than computational work. A multi-pronged optimization strategy can close this gap. A fused murmurate kernel executing all four phases in a single GPU launch would eliminate $\sim\!8$ intermediate HBM round-trips per round, providing $3$--$5\times$ speedup. Increasing batch size to $\geq 8$ amortizes launch overhead, providing $4$--$8\times$ additional effective speedup. Flash-style tiling over the $N$ dimension while keeping the slot pool in SRAM yields another $2$--$3\times$. Combined, these optimizations theoretically enable $24$--$120\times$ speedup, bringing wall-clock time within $1$--$2\times$ of FlashAttention at $N=8{,}192$ while using $9\times$ fewer total FLOPs.

\subsection{The Memory Argument for Decoding}

During autoregressive generation, standard attention must store a KV cache of size $N \cdot D$ per layer. For $N = 100{,}000$ tokens and $D = 4{,}096$, this is 1.6 GB per layer — 19.2 GB for a 12-layer model. Murmurative attention stores a fixed slot pool of size $M \cdot D_h \cdot H = M \cdot D = 65{,}536$ floats per layer, regardless of sequence length. For $N = 100{,}000$ tokens, this represents a $1{,}500\times$ memory reduction in the attention cache, enabling deployment on consumer hardware for arbitrarily long conversations.

% ============================================================
\section{Conclusion}
% ============================================================

Murmurative Attention demonstrates that sub-quadratic attention can match the representational quality of full pairwise attention while using dramatically fewer FLOPs. The core insight — replacing $N^2$ token-token comparisons with $N \cdot M$ token-slot interactions, where $M$ is a small constant (256) — is simple, principled, and empirically validated across sequence lengths from 512 to 8{,}192 tokens.

The remaining gap is engineering, not algorithmics. Our CUDA implementation proves the FLOP advantage; a fused kernel, practical batch sizes, and flash-style tiling should bring wall-clock performance to within $1$--$2\times$ of FlashAttention while preserving the $O(N \cdot M)$ asymptotic advantage. The slot-based attention paradigm opens several research directions: hierarchical slot pools for multi-scale reasoning, learned slot initialization strategies, and extending the murmuration metaphor to cross-attention in encoder-decoder architectures.

We release the full implementation, CUDA kernels, and benchmark suite at \url{https://github.com/beaglabs/murmurative-attention}. If attention is the right mechanism but the wrong complexity class, then slot-mediated attention — murmuration — is the correction.

% ============================================================
\bibliographystyle{unsrt}
\bibliography{references}
% ============================================================

\end{document}